\scanpage{156}\mspage{11}
C'était un complexe de pro-groupes algébriques sur $k$, affines à l'exception de $J_0$ si je me rappelle bien, (\deleted{\unsure{qui}}$J_0$ avait comme partie abélienne la partie abélienne de $\underline{\mathrm{Alb}}_{X/k}$
\margin{\begin{tabular}{@{}|p{\dimexpr\marginparwidth-0.8em\relax}@{}}
\illegible{} de $X$, i.e on avait un \deleted{\unsure{f}}hom $X\longrightarrow\underline{H}_0$ universel pour les hom de $X$ dans des groupes loc. proalgébriques. Pour $X$ connexe, $\underline{H}_0$ est une ext. de $\mathbb{Z}$ par un groupe proalg. connexe.
\end{tabular}}
\tikz[baseline=0pt]{\draw (0,.95ex)--(-.25em,.45ex)--(-.06em,.55ex)--(-.22em,.10ex)--(-.07em,.19ex)--(-.15em,-.14ex);}la jacobienne généralisée ordinaire). Sa construction, inspirée du complexe résiduel, passait par des jacobiennes généralisées (en un sens \inserted{cohomologique} convenable) des localisés $\operatorname{Spec}\underline{O}_{X,x}$ de $X$ en ses différents points. Il est possible qu'en un premier temps, je me limitais au cas $X$ lisse. Le rôle cohomologique de ce complexe était celui d'un complexe d'\uline{homologie}
\[
\mathrm{H}^{i}(X,G_X)\simeq\mathbb{E}\mathrm{xt}^{i}(J_{*X/k},G)\tag*{\makebox[0pt][r]{\makebox[\linewidth][l]{\unsure{\ensuremath{\alpha})|a)|(*)}}}}
\]
mais pour quels coefficients ? Je pense que je devais \deleted{\unsure{me limiter,|me limiter} plus ou moins à des} \inserted{prendre des groupes alg. comm. $G$ quelconques mais travailler en top. de Zariski (malédiction !)\unsure{, --|,, \ensuremath{\top}|. --}} \deleted{$G$ connexes (brr !), ou \unsure{plutôt|plu tôt}} \inserted{même} dans le cas $G$ discret, je prenais néanmoins des $\mathrm{H}^{i}$ zariskiens -- cela donne des groupes de cohomologie un peu bêta, évidemment. Je réalisais qu'il faudrait travailler en coh etale ultérieurement, et que la construction\inserted{\unsure{$\times$|$x$}} des $J_{\deleted{\unsure{*}}iX/k}$ en serait évidemment modifiée. Quant à la \unsure{sig-}%
% archive sheet number: 156
